% Owners: Alejandro drafted 3.1 to 3.5 (items 2a, 2b, 2c, 1d, 3c; app/, src/gates/, src/vlm/prompts.py).
% Austin owns 3.6 (inference and telemetry, items 1a and 1b; src/vlm/, src/energy/) and reviews
% the architecture as a whole (4e plan). Every number below is traced in a comment to its source.
\section{Method}
\label{sec:method}

% Pipeline figure (Alejandro). Stages match app/main.py, src/gates/change.py,
% src/vlm/prompts.py and app/state.py.
\begin{figure*}[t]
\centering
\resizebox{\textwidth}{!}{%
\begin{tikzpicture}[
  font=\footnotesize,
  box/.style={draw, rounded corners=2pt, align=center, minimum height=11mm, inner sep=3pt, fill=#1},
  box/.default=white,
  arr/.style={-{Stealth[length=2mm]}, thick},
  skip/.style={-{Stealth[length=1.6mm]}, gray, dashed},
  lab/.style={font=\scriptsize, gray!40!black, align=center},
  node distance=8mm
]
\node[box=gray!10] (cam) {Webcam or\\real-time replay};
\node[box, right=of cam] (slot) {Latest-frame\\slot (one frame)};
\node[box, right=of slot] (tick) {Check every $\Delta$\\$t_k=t_0+k\Delta$};
\node[box=orange!12, right=of tick] (gate) {Change gate\\score $s$ vs.\ reference\\$\tau{=}0.3$, $k{=}3$, $T_{\min}{=}2$\,s};
\node[box=blue!8, right=of gate] (vlm) {llama-server\\SmolVLM2-500M Q8\_0\\JSON-schema grammar};
\node[box=orange!12, right=of vlm] (val) {Answer\\validation};
\node[box=orange!12, right=of val] (scene) {Gesture\\scene state};

\draw[arr] (cam) -- (slot);
\draw[arr] (slot) -- (tick);
\draw[arr] (tick) -- node[lab, above]{idle} (gate);
\draw[arr] (gate) -- node[lab, above]{send} (vlm);
\draw[arr] (vlm) -- (val);
\draw[arr] (val) -- (scene);

% Skips and feedback.
\node[lab, below=7mm of tick] (busy) {busy: skip\\(nothing queued)};
\draw[skip] (tick) -- (busy);
\node[lab, below=7mm of gate] (gskip) {no change, persisting\\or cooldown: skip};
\draw[skip] (gate) -- (gskip);
\draw[skip] (vlm.north) to[bend right=22] node[lab, above]{sent frame becomes the reference} (gate.north);
\node[lab, below=7mm of val] (unk) {invalid: unknown\\(never \texttt{none})};
\draw[skip] (val) -- (unk);
\node[lab, below=7mm of scene] (ev) {onset / offset\\events};
\draw[arr] (scene) -- (ev);

% Logs.
\node[lab, below=15mm of vlm, text=black] (logs) {Run directory: requests, drops, predictions, power, manifest (versioned file contract)};
\end{tikzpicture}}
\caption{The assistant's pipeline. Orange stages are the parts described in Secs.~\ref{sec:gate} to \ref{sec:scene}. The fixed-interval baseline is the same pipeline with the gate removed, so it sends every check at which no request is in flight. Every skipped check is logged with its reason.}
\label{fig:pipeline}
\end{figure*}


The assistant watches one camera and reports when a person starts or stops showing a thumbs-up or a peace sign.
Figure~\ref{fig:pipeline} shows the pipeline.
A capture thread keeps only the newest frame, a scheduler decides at regular checks whether that frame goes to the vision-language model (VLM), the model's answer is constrained to a fixed format and validated, and a scene state turns the stream of validated answers into gesture events.
The fixed-interval baseline and the gated policy share every stage except the gate, so any difference between them comes from which frames are sent.

\subsection{Capture and scheduling}
\label{sec:schedule}
% Source: app/main.py module docstring and _main_loop.
Capture runs on its own thread and overwrites a single slot with each new frame, so slow inference never stalls capture and no backlog of stale frames builds up.
The scheduler wakes at checks $t_k = t_0 + k\Delta$ on the monotonic clock, so checks do not drift when one runs late.
At each check it takes the newest frame and applies three rules in order.
If a request is still in flight, the check is skipped and logged as busy; nothing is queued.
If an earlier request timed out, the scheduler waits until the server reports no busy slot before sending again.
Otherwise the frame is sent under the fixed-interval policy, or passed to the change gate under the gated policy.
At most one request is in flight at any time.
Every skipped check is written to the run's drop log with its reason, so the logs account for every check.

\subsection{Change gate}
\label{sec:gate}
% Source: src/gates/change.py (GateConfig defaults and ChangeGate.decide/acknowledge).
The gate compares the candidate frame with a \emph{reference}: the last frame actually sent to the VLM.
Comparing against the last sent frame, instead of the previous frame, lets a change that builds up slowly keep growing until it triggers.

\paragraph{Change score.}
Each frame is converted to grayscale, resized to $160\times120$, blurred with a $5\times5$ Gaussian kernel, and its mean brightness is subtracted, so a uniform lighting shift does not count as change.
For a preprocessed frame $f$ and reference $r$, a pixel $p$ is changed when $|f(p)-r(p)| > \delta$, with $\delta=25$ on the 0 to 255 brightness scale.
The image is divided into a $4\times4$ grid of cells $C_j$, and the score is the changed fraction of the most-changed cell,
\begin{equation}
  s = \max_j \frac{1}{|C_j|}\sum_{p\in C_j} \mathbb{1}\big[\,|f(p)-r(p)|>\delta\,\big].
  \label{eq:score}
\end{equation}
Taking the maximum over cells keeps a hand raised in one part of the frame from being diluted by the static remainder.
The whole-frame changed fraction is logged beside $s$ for diagnosis.

\paragraph{Decision.}
A change starts when $s \ge \tau$ and must stay above $\tau$ for $k$ consecutive checks before it triggers a call; a score below $\tau$ resets the count.
After a trigger, further triggers are held back for a cooldown $T_{\min}$.
In order, the gate (1) triggers when it has no reference yet, (2) skips when $s<\tau$, (3) skips while the change has lasted fewer than $k$ checks, (4) skips within $T_{\min}$ of the last trigger, and (5) otherwise triggers.
When a triggered frame is sent, it becomes the new reference, the cooldown clock restarts and the persistence count resets.
The gate is consulted only when no request is in flight, so every trigger is sent, and a busy check neither advances nor resets the count.
The implementation also supports a periodic forced refresh and a lower release threshold (hysteresis); both are off by default.

\paragraph{Settings.}
% Source: GateConfig defaults; README "Choosing the gate settings" for the tuning evidence.
We use $\tau=0.3$, $k=3$ and $T_{\min}=2$\,s with checks every $\Delta=0.5$\,s, so a change must be present at three checks spanning 1\,s before it is sent.
These values were chosen on development clips.
On \texttt{real\_thumbs\_1}, ordinary movement scores 0.02 (median) to 0.40 (99th percentile), while a thumbs-up peaks at 0.68 to 0.75.
An earlier setting ($\tau=0.05$ without persistence) fired on head movement and made more calls than fixed 2\,s sampling.
With $\tau=0.3$ and $k=3$, a replay that models each call as 2.5\,s of busy time caught both gestures on \texttt{real\_thumbs\_1} with 27 calls against 147 for fixed 2\,s sampling, 82\% fewer.
A threshold of 0.4 cut calls further but left less margin under the gesture scores, and hysteresis only added calls on these clips.
Because $k$ counts checks, these settings are tied to $\Delta=0.5$\,s and must be retuned for another check interval.

\subsection{Validated answer format}
\label{sec:answers}
% Source: src/vlm/prompts.py (GESTURE_PROMPT, GESTURE_SCHEMA, parse_gesture_answer, ANSWER_ERRORS);
% configs/inference.json (max_output_tokens 64).
Every request asks the same question about one frame: \emph{``Describe what the person's hands are doing in a few words, then say whether the person is making a thumbs up, a peace sign, or neither (none).''}
The request carries a JSON schema that llama-server compiles into a decoding grammar, so a compliant server can only return an object with two fields in a fixed order: \texttt{hands}, a free description of at most 80 characters, then \texttt{gesture}, one of \texttt{thumbs\_up}, \texttt{peace\_sign} or \texttt{none}.
The \texttt{hands} field comes first because SmolVLM2 asked for the gesture alone answered \texttt{thumbs\_up} on every frame; writing a short description first gives the model a step in which to look before it commits to a label.
Only \texttt{gesture} is used as the observation.

The application validates every answer exactly, without normalizing or guessing, because the constraint can fail: a server may ignore the schema, or the 64-token output bound may cut an answer short.
An answer that fails validation is recorded with one of seven reasons (truncated, not JSON, not an object, missing field, unexpected fields, non-text description, gesture outside the vocabulary) and is treated as \emph{unknown}.
An unknown answer is never read as \texttt{none}: \texttt{none} is a positive statement that neither gesture is visible, while an unknown answer carries no information.
Transport status (whether the server answered) is logged separately from answer validity.

\subsection{Gesture scene state}
\label{sec:scene}
% Source: app/state.py (GestureScene.update, Evidence.available_source_time_s).
The scene state holds one confirmed gesture, starting at \texttt{none}, and updates once per request.
A validated answer equal to the confirmed gesture changes nothing, so a held gesture is reported once.
A different validated answer confirms a change after $c$ consecutive answers agree ($c=1$ by default): leaving a gesture emits an \emph{offset}, entering one emits an \emph{onset}, and a direct switch between the two gestures emits both from the same answer.
Returning to \texttt{none} rearms the state, so the same gesture shown again is a new onset.
Unknown answers (failed requests or invalid answers) change nothing: the confirmed gesture stays, and a change awaiting confirmation neither advances nor resets.
The state also counts unknown answers since the last valid one, which the display shows as staleness.

Each event keeps the request it came from, with the frame's source time and the time the answer became available.
The event time used for scoring is the availability time mapped onto the source timeline,
$t_{\text{avail}} = t_{\text{src}} + (t_{\text{resp}} - t_{\text{cap}})$,
where $t_{\text{src}}$ is the frame's media time, $t_{\text{cap}}$ its host capture time and $t_{\text{resp}}$ the host time the answer arrived.
Response delay therefore includes the time the frame waited to be sent and the inference time, and does not credit an answer with the moment its frame was captured.
Both policies feed the same scene state, so fixed and gated runs share one event derivation.

\subsection{Run records}
\label{sec:records}
% Source: app/main.py module docstring; docs/data-contracts.md (item 2d); src/contracts/.
Each run writes a new directory with a request log (one line per call, with frame identity, source time, request and response times, transport status, validated gesture or answer error), a drop log (every skipped check and every frame the source dropped, with its reason and gate score), the derived gesture events, power samples and a manifest of the configuration and platform.
The formats follow the project's versioned file contract, and a validator checks every run directory, so the evaluation tools read files and never application internals.
Unavailable power, token or timing values are recorded as missing, never as zero.

\subsection{Inference and telemetry}
\label{sec:inference}
% Owner: Austin (items 1a, 1b). Alejandro: facts below are from configs/inference.json and
% docs/1b-battery-qualification.md; Austin turns them into prose and checks them.
\TODO{Inference and telemetry prose: Austin.}
\begin{notes}{Austin}
\begin{itemize}\setlength\itemsep{0pt}
\item Model: SmolVLM2-500M-Video-Instruct, Q8\_0 GGUF with its Q8\_0 projector (\texttt{ggml-org/SmolVLM2-500M-Video-Instruct-GGUF}, pinned revision and SHA-256 in \texttt{configs/inference.json}).
\item Runtime: one persistent llama-server from llama.cpp~\cite{llamacpp} revision 620d222, CPU only, 4 threads, context 4096, output bound 64 tokens, 120\,s read timeout, localhost port 8080.
\item Client: \texttt{src/vlm/client.py}; failed requests are logged, never retried silently. Frames sent as JPEG at quality 85 (\texttt{app/main.py}).
\item Team-laptop smoke check (1a, \texttt{docs/validation-1a-1b.md}): mean request latency 2.48\,s under WSL2 on CPU.
\item Telemetry (1b): whole-laptop battery energy from Windows \texttt{BatteryReport} read from WSL2; WSL sysfs scale is unqualified and RAPL is absent, so package energy is unavailable. Details in Sec.~\ref{sec:setup}.
\end{itemize}
\end{notes}
